\chapter{Independent prime divisors}\label{ch:oc-divisors}

This chapter treats the prime factors of a typical integer as almost independent, from the Tur\'an--Kubilius inequality to the Erd\H os--Kac theorem.

\section{The Tur\'an--Kubilius inequality}

The Tur\'an--Kubilius inequality bounds the mean square of an additive function by a sum over primes. It is the second-moment form of the idea that the events ``$p$ divides $n$'' are nearly independent.

\section{The Erd\H os--Kac theorem}

Let $\omega(n)$ be the number of distinct prime factors of $n$. For $n$ drawn uniformly from $\{1,\ldots,\lfloor x\rfloor\}$,
\[
\frac{\omega(n)-\log\log x}{\sqrt{\log\log x}}
\]
tends in distribution to the standard Gaussian law. That is the Erd\H os--Kac theorem.

\section{The Kubilius model}

The Kubilius model replaces those events by independent Bernoulli random variables of probabilities $1/p$. The chapter's aim is to say for which additive functions the model and the integers give the same limit law.
